\section{Appendices}
\label{sec:appendices}

\subsection*{A. Residual-risk acknowledgment for the pure-PQ core promise}

\pqcrypt{}'s core promise (suites \code{0x01} and \code{0x02}) is post-quantum-only. If ML-KEM-1024 is later structurally broken (a sub-exponential Module-LWE attack), every \code{0x01}/\code{0x02} envelope loses confidentiality with no classical fallback. This is a 30--50 year risk accepted in exchange for the unambiguous PQ-only promise. The mitigation path is now operational rather than aspirational:

\begin{enumerate}
\item Suite byte \code{0x06} (PQCK-00021, v1.2; \cref{sec:hybrid}) is the opt-in hybrid suite (X25519 + ML-KEM-1024 + AES-256-GCM + UtC). Callers who want belt-and-suspenders coverage select \code{0x06} at \code{encrypt} time. The wire byte at offset 7 records which promise the envelope makes.
\item Long-lived ciphertexts written under suites \code{0x01}/\code{0x02} can be re-encrypted under \code{0x06} (or a future suite) within a defined window. \pqcrypt{} v1.x is not an archive format; long-lived storage should plan re-encryption.
\item Suite \code{0x01}/\code{0x02} decoders remain available indefinitely (suite bytes are immutable post-allocation, matching the variant-byte discipline of \pqid{} and \pqhash{}).
\end{enumerate}

We do \emph{not} mitigate by hedging suites \code{0x01}/\code{0x02} themselves. The pure-PQ promise of those suites is preserved verbatim from v1.0; the hybrid escape valve sits in a separate suite-byte allocation. This is the design choice that lets the same wire format carry two distinct guarantees without diluting either.

\subsection*{B. Variant byte allocation across tessera}

\begin{center}
\small
\begin{tabular}{@{}llll@{}}
\toprule
\textbf{Byte} & \textbf{Owner} & \textbf{Meaning} & \textbf{Status} \\
\midrule
\code{0x00} & --- & Permanently reserved & MUST reject \\
\code{0x01} & \pqid{} & Standard generator & assigned \\
\code{0x02} & \pqid{} & HLC generator & assigned \\
\code{0x03} & \pqhash{} & SHAKE-256 / 64 B tail & assigned \\
\code{0x04} & \pqcrypt{} & Encryption envelope v1 & assigned (this document) \\
\code{0x05}--\code{0xFE} & --- & Reserved-future & MUST reject \\
\code{0xFF} & --- & Permanently reserved & MUST reject \\
\bottomrule
\end{tabular}
\end{center}

Cross-rejection is normative: \pqid{} and \pqhash{} decoders MUST reject \code{0x04}; \pqcrypt{} decoders MUST reject \code{0x01}, \code{0x02}, \code{0x03}, and every byte in \code{0x05}--\code{0xFF} in v1. Each project's \texttt{invalid.json} conformance file pins the cross-rejection.

\subsection*{C. Suite byte allocation}

\begin{center}
\small
% \resizebox: the natural table width exceeds \textwidth; without it the Status
% column renders off the page edge (pre-existing clip, fixed TSEAL-00048).
\resizebox{\textwidth}{!}{%
\begin{tabular}{@{}lll@{}}
\toprule
\textbf{Suite} & \textbf{Algorithms} & \textbf{Status} \\
\midrule
\code{0x01} & ML-KEM-1024 + AES-256-GCM + UtC commit + HKDF-SHA3-256                & v1.0 mandatory (pure PQ) \\
\code{0x02} & ML-KEM-1024 + ChaCha20-Poly1305 + UtC commit + HKDF-SHA3-256          & v1.0 optional (pure PQ) \\
\code{0x06} & X25519 + ML-KEM-1024 + AES-256-GCM + UtC commit + HKDF-SHA3-256       & v1.2 optional (hybrid) \\
\code{0x07} & RESERVED: ML-KEM-1024+HQC hybrid or HQC-only (\cref{sec:suites})       & reserved; MUST reject \\
\code{0x00}, \code{0x03}--\code{0x05}, \code{0x08}--\code{0xFF} & Reserved-future   & MUST reject \\
\bottomrule
\end{tabular}}
\end{center}

Earmarked future suites (not normative in v1.x):
\begin{itemize}
\item \code{0x03}: ML-KEM-1024 + AEGIS-128L + UtC commit. Evaluated and deferred in PQCK-00020 (v1.2, 2026-05-21): AEGIS libraries not yet polyglot-ready. The byte remains reserved-future; revisit when libcrux, OpenSSL, and Go stdlib all ship verified AEGIS implementations.
\item \code{0x07}: \textbf{explicitly RESERVED} (TSEAL-00048) for the HQC contingency --- an ML-KEM-1024+HQC hybrid or HQC-only successor suite \citep{nisthqc2025}; allocated upon NIST HQC FIPS finalization (draft slipped; final $\sim$2027); activation trigger = Module-LWE confidence loss or HQC standardization. See \cref{sec:suites} for the full roadmap paragraph.
\item Hardened-implementation variants (masked ML-KEM, side-channel-protected X25519) would each take a new suite byte rather than modify \code{0x01}/\code{0x02}/\code{0x06}.
\end{itemize}

\subsection*{D. Test-vector schema}

Each KAT file is a JSON array of records. Common fields:

\begin{lstlisting}[language=Pseudocode]
{
  "name":         "human-readable id, unique within the file",
  "op":           "encrypt | decrypt | peek_header | commit | ...",
  "input_hex":    "raw bytes for decode-from-binary ops",
  "input":        { ... op-specific structured input },
  "expected":     { ... op-specific structured output },
  "should_fail":  false,
  "error_kind":   "DECRYPT_FAILED | UNKNOWN_VARIANT | UNKNOWN_SUITE | INVALID_LEN"
}
\end{lstlisting}

Full schema in \texttt{tessera-seal/conformance/README.md}.

\subsection*{E. Acknowledgments}

\pqid{} and \pqhash{} for the polyglot-spec-first discipline this whitepaper inherits. Bellare and Hoang for the UtC construction \citep{bellare2022commitment}. The OpenSSL maintainers for the ML-KEM-1024 \texttt{EVP\_KEM} integration shipped in OpenSSL 3.5; the RustCrypto \texttt{ml-kem} authors for the pinned Rust source; the Go core team for landing \texttt{crypto/mlkem} in the 1.26 stdlib; the \texttt{kyber-py} maintainers (NCC Group audit, Q2 2024); and the Cryspen team for \texttt{mlkem-native} \citep{cryspen2025mlkem}, evaluated during the v1.0 drafts (the shipped C reference uses OpenSSL instead; \cref{sec:mlkem}). The X25519 sources used by hybrid suite \code{0x06}: \texttt{x25519-dalek} v2 (Rust), Go stdlib \texttt{crypto/ecdh}, and \texttt{cryptography} (Python). The X-Wing draft authors \citep{xwing-draft} for the combiner design idiom. The KyberSlash \citep{kyberslash}, PQShield \citep{pqshield2024}, SPA-CCA \citep{guo2024spa}, and GoFetch \citep{gofetch} disclosure teams, whose findings shape the side-channel disclaimer.

\subsection*{F. Release timeline (v1.0 -- v1.3.x)}
\label{sec:appendix:timeline}

\begin{center}
\small
\begin{tabularx}{\textwidth}{@{}lllX@{}}
\toprule
\textbf{Release} & \textbf{Date} & \textbf{Tickets} & \textbf{Headline} \\
\midrule
v1.0 alpha   & 2026-05-20 & PQCK-00001--00010 & Spec + Python/C/Go/Rust references + plugin protocol + cross-language conformance harness. Suites \code{0x01}/\code{0x02} only. \\
v1.1         & 2026-05-21 & PQCK-00011--00019 & C++ / TypeScript / Java / CPython ext ports; \texttt{.tseal-key-v1.1} keyfile format (embeds pubkey); CSPRNG path rewire; AEGIS-128L evaluation (deferred). C switched from \texttt{mlkem-native} to OpenSSL 3.5 \texttt{EVP\_KEM}; Java switched from BouncyCastle 1.79+ to SunJCE JEP 452; Go switched from \texttt{filippo.io/mlkem} to stdlib \texttt{crypto/mlkem}. \\
v1.2         & 2026-05-21 & PQCK-00020--00026 & \textbf{Hybrid suite \code{0x06}} (Python only); reference CLI binary; GitHub Actions matrix; bench sweep; internal security audit doc; AEGIS deferral recorded. \\
v1.3 partial & 2026-05-21 & PQCK-00027--00029 & Hybrid suite \code{0x06} ports to Rust (PQCK-00028) and Go (PQCK-00029). AES-NI/ARMv8 intrinsics deferral for Rust (PQCK-00027). \\
v1.3 hardening & 2026-06-08--10 & TSEAL-00035--00045 & BadKey suite/key consistency completed across all eight ports (TSEAL-00035--00038); CLI suite-flag hard-error (TSEAL-00039); keyfile KDF block, fail-closed (TSEAL-00040); KAT expansion + append-only vector freeze; suite-\code{0x06} rejection tests in the five non-hybrid ports (TSEAL-00044); whitepaper truthfulness errata (TSEAL-00045, this revision). \\
\bottomrule
\end{tabularx}
\end{center}

v1.4 backlog (not normative in this document): the hybrid suite \code{0x06} port to C landed (TSEAL-00030, 2026-06-13 --- a 4-way byte-identical KAT contract with Python/Rust/Go), and library-level hybrid then landed in C++ (TSEAL-00031), TypeScript (TSEAL-00032), and the CPython extension (TSEAL-00034) on the same day, each verified by an in-suite test that decodes the frozen \texttt{hybrid\_*.json} vectors byte-identically without a tessera-seal conformance runner; Java (TSEAL-00033) is the sole remaining hybrid port, pending a JDK 24+ build/verify environment. Also outstanding: adding tessera-seal conformance runners for C++, TypeScript, Java, and the CPython extension so their (already library-level) coverage is harness-enforced; and reopening the Rust AES-NI/ARMv8 intrinsics work deferred in PQCK-00027. The automated CI timing-histogram regression (\texttt{security-audit-v1.2.md} open item \#1) landed in TSEAL-00048 (an earlier revision listed it as PQCK-00035, a ticket number since consumed by the BadKey suite/key-consistency work released 2026-06-08). Ticket-prefix note: PQCK-\emph{NNNNN} identifiers are historical (pre-rename); the same numbering continues as TSEAL-\emph{NNNNN}, and \texttt{tessera-seal/CLAUDE.md} is the authoritative ticket registry where this document and it disagree.

\subsection*{G. Author and document provenance}
\label{sec:appendix:provenance}

\textbf{Author:} Bahram Samani, \texttt{tessera} monorepo maintainer.

This document is built from LaTeX sources under \texttt{tessera-seal/latex/} (\texttt{main.tex} + \texttt{secNN\_*.tex}); the build chain is \texttt{pdflatex $\rightarrow$ bibtex $\rightarrow$ pdflatex $\times$ 2} per the Makefile. The artifact at \texttt{tessera-seal/latex/TesseraSeal\_Whitepaper.pdf} is regenerated and committed alongside any normative change.

\textbf{Authority order} (when sources disagree): this whitepaper, then the Python reference at \texttt{tessera-seal/python/tessera\_seal.py}, then the canonical C ABI at \texttt{tessera-seal/c/tessera\_seal.h}, then every other language. A cross-language disagreement is never resolved by changing the spec to match an outlier; the outlier is fixed.

\textbf{Threat-model checklist}: the internal pre-external-audit walk-through is at \texttt{tessera-seal/docs/security-audit-v1.2.md} (PQCK-00025). It enumerated six open items for an external auditor; of these, the CI timing histogram (automated in TSEAL-00048) and the hybrid C port (TSEAL-00030, 2026-06-13) have since landed, leaving four (FIPS 140-3 path, synthetic-key timing-pad invariant testing, plugin protocol fuzz harness, memory-zeroization audit).

\textbf{Side-channel posture}: this document and the companion \texttt{tessera-seal/docs/side-channel-disclaimer.md} are normative. v1.x mitigates software-timing side channels only; not power, EM, frequency-scaling, or prefetcher attacks.

\textbf{Conformance gate}: \texttt{tools/conformance/run.sh} backs the CI gate --- \texttt{.github/workflows/conformance.yml} runs on every pull request and push, and branch protection on \texttt{main} requires its cross-language job (the 4-language alpha step and the TesseraSeal timing-histogram step block; the full 8-language step is informational). Latest run (2026-06-17, Apple M4 Max alpha quick loop): 1192 PASS, 0 FAIL across \pqid{} + \pqhash{} + \pqcrypt{}, four alpha languages ($298 + 298 + 298 + 298$ --- includes the append-only \texttt{kyberslash\_regression.json} record; \cref{sec:conformance}).
